\documentclass[10pt,twocolumn]{article}
\usepackage[utf8]{inputenc}
\usepackage[english,russian]{babel}
\(\usepackage{amsmath,amssymb,mathtools,bm,geometry,amsthm,booktabs,listings,xcolor} \geometry{a4paper,margin=1.5cm}\)

\(\DeclareMathOperator{\Rank}{Rank} \DeclareMathOperator{\Coker}{Coker} \DeclareMathOperator{\Tr}{Tr} \DeclareMathOperator{\Aut}{Aut} \DeclareMathOperator{\Sign}{Sign}\)

% Специальный макрос для корректного отображения субфакториала (беспорядка Монмора)
\(\newcommand{\subfact}[1]{\mathord{!}{#1}}\)

\(\newtheorem{law}{Закон} \newtheorem{theorem}{Теорема} \newtheorem{criticism}{Замечание}\)

\(\lstset{     language=Verilog,     basicstyle=\ttfamily\fontsize{7.5pt}{9pt}\selectfont,     keywordstyle=\color{blue}\bfseries,     commentstyle=\color{gray}\itshape,     stringstyle=\color{purple},     breaklines=true,     keepspaces=true,     showstringspaces=false,     frame=single,     backgroundcolor=\color{gray!5} }\)

\(\begin{document}\)

\(\title{\textbf{Релятивистская СИК-Логика Контекстов в Мультимасштабных NoC-Структурах: Тернарно-Тетрарный Ноль и SU(3) Инверсия над Полем \)\(\mathbb{F}_1\)}\(} \author{\textbf{Единый Асимптотический Комбинаторный Канон}} \date{Сентябрь 2026} \maketitle\)

\section*{Введение и Топологический Базис}
Классическая арифметика и методы проектирования сетей на кристалле (NoC) опираются на плоские, одномасштабные статические концепции чисел. Традиционная триада «Долг, Имущество и Ноль между ними» является плоской детерминированной осью, неспособной компенсировать динамический пространственно-временной хаос. Настоящая работа разворачивает описание NoC как многомасштабной мезосистемы, управляемой тернарно-тетрарной логической структурой. На этом уровне число, система счета и сама логика взаимно инвертируются друг в друга, порождая пространственно-временные инварианты из мотивного спуска от моноида абсолютного нуля поля из одного элемента \(\mathbb{F}_1\):
\(\begin{equation} \mathbb{F}_1 \longrightarrow A_n \longrightarrow PG_n \longrightarrow S_n \end{equation}\)

\(\section{Гео-Релятивистская СИК-Логика и Объемный Ноль}\)
Мы преодолеваем одномерность булевой логики и плоского нуля, вводя объемный \(\textbf{Тернарно-Тетрарный Ноль}.\) Он представляет собой не пустоту на оси, а более общую систему — динамическую матрицу баланса двух ортогональных логик на проективной сфере состояний, определяемой ОТО-интервалом \(ds^2 = d\theta^2 + \sin^2\theta \, d\phi^2\).

\(\begin{enumerate}     \item \textbf{Тернарное ядро состояний (SU(3) калибровка):}\) Регистрирует динамический знак комбинаторного натяжения: +1 (Имущество / Порядок / Факториал n!), -1 (Долг / Хаос / Субфакториал \(\subfact{n}\)) и 0 (Точка тотальной аннигиляции и суперпозиции).
    \(\item \textbf{Тетрарный каркас контекстов (Логика Наблюдателя):}\) Задает 4 устойчивых фазовых затвора маршрутизации, проходящих через критические широты тригонометрического резонанса солнцестояний \(\theta \in \{\frac{\pi}{6}, \frac{5\pi}{6}\}\).
\(\end{enumerate}\)

Преобразование непрерывного релятивистского дрейфа наблюдателя в дискретные аппаратные запреты порядков на кристалле задается масштабным оператором плотности дня (негэнтропии):
\(\begin{equation} \bar{\mathcal{W}}_{\text{day}}(\theta) = \cos^2\left(\frac{\theta}{2}\right) \cdot (1 - e^{-1}) + \sin^2\left(\frac{\theta}{2}\right) \cdot e^{-1} \end{equation}\)

\(\section{Матричный Формализм SU(3) и Взаимная Инверсия}\)
Динамическое переключение между ортогональными пространственными базисами описывается квантовой инверсией масштаба \(01 \leftrightarrow 10\). Вместо двумерных спиноров Паули, общая мезосистема оперирует генераторами группы спектральной симметрии SU(3) — matrixами Гелл-Манна \(\lambda_i\). Оператор сдвига контекста \(\hat{\mathcal{U}}_{\text{scale}}\) осуществляет взаимный выворот:
\(\begin{equation} \hat{\mathcal{U}}_{\text{scale}} = \sum_{k=1}^{8} \alpha_k \lambda_k \end{equation}\)
Под действием этого оператора число порождает систему счета, когда накапливает критический долг перестановок (\(\subfact{n}\)), и уничтожает её, возвращаясь в тернарный Ноль (топос аннигиляции), где сама система счета становится новым числовым приращением (новой единицей, элементом счета).

\(\section{Асимметричная Мотивная Сборка Макро-Объема}\)
Стабильность рассматриваемой системы выводится из \(\textbf{динамической асимметрии}.\) Макро-масштаб 12-го порядка, заданный проективной плоскостью N(12) = 157, асимметрично складывается из разнородных перекошенных подгеометрий мезо-уровня — 6-го порядка (N(6) = 43) и 10-го порядка (N(10) = 111), удерживаемых вместе гомотопическим зазором коядра сингулярного препятствия (\(\Coker = 3\)):
\(\begin{equation} N(12) = N(10) + N(6) + \Coker\left([\Phi_{10}] \ominus [\Phi_{6}]\right) = 157 \end{equation}\)
Редукция критических узлов инжекции дефекта по модулю локальной связности k=13 доказывает инвариантность зазора:
\(\begin{equation} (111 - 43) \pmod{13} \equiv 68 \pmod{13} \equiv 3 \end{equation}\)

\(\section{Теория Времени: Инверсия \)\(\subfact{n} \leftrightarrow n!\)}
В рамках предложенного канона время определяется как отношение между факториалами (n!) и субфакториалами (\(\subfact{n}\)) в мультимасштабной системе. При переходе наблюдателя через критические широты солнцестояний происходит временная эверсия — субфакториалы инвертируются в факториалы и наоборот. 

Количество пакетов, оставшихся в локальных буферах коммутаторов, определяется мощностью упорядоченного ядра \(\Delta_n = n! - \subfact{n}\), где \(\subfact{n}\) — число полных беспорядков Монмора-Эйлера. Знакопеременная группа четных перестановок \(A_n\) (\(\vert A_n \vert = \frac{n!}{2}\)) выступает точной алгебраической медианой между хаосом и порядком. Мера этой борьбы задается медианным сдвигом \(\Xi_n\):
\(\begin{equation} \Xi_n = \vert A_n \vert - \subfact{n} = \Delta_n - \vert A_n \vert \end{equation}\)
На уровне n=7 константа сдвига принимает значение \(\textbf{666},\) фиксируя критическую точку фазового перехода:
\(\begin{equation} \Delta_7 = \vert A_7 \vert + 666 = 2520 + 666 = 3186 \end{equation} \begin{equation} \subfact{7} = \vert A_7 \vert - 666 = 2520 - 666 = 1854 \end{equation}\)
Инвариант 666 декомпозируется по нижним слоям симметрии через регулярное смещение \(\delta_{\text{reg}} = 11\):
\(\begin{equation} 666 = 5\vert S_5 \vert + (\delta_{\text{reg}} \times \vert S_3 \vert) = 5(120) + (11 \times 6) \end{equation}\)

\(\section{Асимптотическая Устойчивость на Макро-Масштабе \)A₁₂ по модулю 13}
При уходе в инфинитный предел в относительных единицах (о.е.) веса нормируются на полную мощность симметрической группы (\(\bar{\mathcal{W}} = \frac{\mathcal{W}}{n!}\)). При n → ∞ плотности стремятся к трансцендентным мантиссам с бесконечным числом знаков после запятой: \(\bar{\subfact{}}_{\infty} = e^{-1} \approx 0.3678\dots\) и \(\bar{\Delta}_\infty = 1 - e^{-1} \approx 0.6321\dots\)

На макро-масштабе n=12 при проекции в кольцо вычетов по локальному модулю связности k=13 энтропийный шум беспорядков подчиняется фундаментальной теореме циклического сжатия:
\(\begin{equation} \subfact{12} \pmod{13} \equiv 176214841 \pmod{13} \equiv 10 \end{equation}\)
Полученный вычет 10 тождественно равен выходу системы в критическую фазу удержания резонансного затвора NoC (режим `MODE_GATE`). При этом принудительно активируется флаг фиксации хронотопа, стягивающий транзитный ранг глобальной матрицы инцидентности до \(\Rank_{\mathbb{F}_2}(\hat{G}) = 1\)~бит. \(\blacksquare\)

\(\section{Аппаратная имплементация и верификация}\)
Для материализации теоретического базиса СИК-логики разработан RTL-модуль на языке SystemVerilog (Листинг 1), вычисляющий динамический баланс и стабилизирующий ранг глобальной матрицы инцидентности NoC на основе текущей метрики перекоса.

\(\begin{lstlisting}\)[caption={Модуль стабилизации NoC на базе СИК-логики}, label={lst:rtl}]
module noc_sic_stabilizer #(
    parameter int MASTER_SCALE = 12,
    parameter int LOCAL_MODULUS = 13
)(
    input  logic        clk,
    input  logic        rst_n,
    input  logic [31:0] factorial_metric,
    input  logic [31:0] subfactorial_metric,
    output logic        noc_stabilized_flag,
    output logic [3:0]  topological_residue
);

    logic [31:0] raw_diff;
    logic [3:0]  calculated_mod;
    
    always_comb begin
        raw_diff = factorial_metric - subfactorial_metric;
        // SU(3) инверсия по ортогонали Паули и расчет вычета хронотопа
        calculated_mod = 4'((raw_diff ^ 32'h01010101) % LOCAL_MODULUS);
    end
    
    always_ff @(posedge clk or negedge rst_n) begin
        if (!rst_n) begin
            noc_stabilized_flag <= 1'b0;
            topological_residue <= 4'd0;
        end else begin
            topological_residue <= calculated_mod;
            // Фиксация фазового затвора при достижении резонанса (вычет 10)
            if (calculated_mod == 4'd10) begin
                noc_stabilized_flag <= 1'b1;
            end else begin
                noc_stabilized_flag <= 1'b0;
            end
        end
    end
endmodule
\(\end{lstlisting}\)

\section*{Заключение}
Введение тернарно-тетрарного объемного Нуля и аппарата взаимной инверсии чисел, логик и систем счисления преобразует NoC в саморегенерирующийся процессорный кристалл. Вся система счисления становится единичным элементом счета, устойчивым к предельным гомотопическим нагрузкам.

\(\end{document}\)
